\subsection{Các phương pháp Học máy hiện có để giải quyết bài toán}
%Ghi rõ và ưu nhược điểm
Dữ liệu mà nhóm nghiên cứu có các tính chất đặc thù như quy mô dữ liệu khổng lồ, tính chất chuỗi thời gian, tồn tại nhiễu cảm biến diện rộng và sự mất cân bằng giữa 2 lớp nghiêm trọng.

\subsubsection{Nhóm Phân Lớp Có Giám Sát (Supervised Classification)}
Do dữ liệu đã có gán nhãn cho biến \texttt{anomaly}( $0$ biểu thị trạng thái bình thường, $1$ biểu thị trạng thái bất thường), họ thuật toán dựa trên cây quyết định (Tree-based Models) và phương pháp học tập hợp (Ensemble Learning) thường thể hiện hiệu suất vượt trội.

\paragraph{XGBoost (eXtreme Gradient Boosting)}\hspace{1cm}

XGBoost tối ưu hóa hàm mục tiêu bằng phương pháp khai triển Taylor bậc hai, cho phép hội tụ nhanh hơn. Cấu trúc cây được phát triển theo từng tầng (Level-wise growth). Mô hình cuối cùng là tổng hợp tuần tự của $K$ cây quyết định:
\begin{equation*}
    \hat{y}_i = \sum_{k=1}^{K} f_k(\mathbf{x}_i), \quad f_k \in \mathcal{F}
\end{equation*}
trong đó $\mathcal{F}$ là không gian của các cây hồi quy CART. Tại mỗi bước lặp $t$, hàm mục tiêu được tối ưu bằng khai triển Taylor bậc hai:
\begin{equation*}
    \mathcal{L}^{(t)} \simeq \sum_{i=1}^{n} \left[ g_i f_t(\mathbf{x}_i) + \frac{1}{2} h_i f_t^2(\mathbf{x}_i) \right] + \Omega(f_t)
\end{equation*}
với $g_i = \partial_{\hat{y}^{(t-1)}} \ell(y_i, \hat{y}^{(t-1)})$ là gradient bậc nhất, $h_i = \partial^2_{\hat{y}^{(t-1)}} \ell(y_i, \hat{y}^{(t-1)})$ là gradient bậc hai (Hessian) của hàm mất mát và $\Omega(f_t)$ là hàm phạt cấu trúc cây:
\begin{equation*}
    \Omega(f) = \gamma T + \frac{1}{2} \lambda \sum_{j=1}^{T} w_j^2 + \alpha \sum_{j=1}^{T} |w_j|
\end{equation*}
trong đó $T$ là số lá, $w_j$ là trọng số tại lá $j$, $\gamma$ phạt số lượng lá, $\lambda$ là hệ số điều chuẩn $L_2$ và $\alpha$ là hệ số điều chuẩn $L_1$.\\

\textbf{Ưu điểm:}
\begin{itemize}
    \item \textbf{Cơ chế điều chuẩn (Regularization):} Tích hợp các thành phần phạt $L_1$ ($\alpha$) và $L_2$ ($\lambda$) trực tiếp vào hàm mục tiêu $\Omega(f)$, giúp hạn chế hiệu quả hiện tượng overfitting khi mô hình tiếp xúc với dữ liệu nhiễu từ các cảm biến đo lường.
    \item \textbf{Xử lý mất cân bằng lớp:} Hỗ trợ điều chỉnh trọng số thông qua siêu tham số \texttt{scale\_pos\_weight}, gán penalty cao hơn cho các trường hợp phân loại sai lớp thiểu số.
    \item \textbf{Khai triển Taylor bậc hai:} Sử dụng cả gradient bậc nhất $g_i$ và Hessian $h_i$ giúp xấp xỉ hàm mất mát chính xác hơn, tạo ra các bước cập nhật tối ưu hơn so với các phương pháp Gradient Boosting truyền thống chỉ dùng gradient bậc nhất.
    \item \textbf{Xử lý giá trị thiếu tự động:} Thuật toán tự học hướng mặc định (default direction) tại mỗi nút chia cho các giá trị bị khuyết, phù hợp với dữ liệu cảm biến thực tế thường có nhiều giá trị bị thiếu.
\end{itemize}

\textbf{Nhược điểm:}
\begin{itemize}
    \item \textbf{Độ phức tạp không gian và thời gian:} Thuật toán phân nhánh tham lam chính xác (Exact Greedy Algorithm) đòi hỏi duyệt qua toàn bộ dữ liệu liên tục. Mặc dù đã có các cải tiến như Histogram-based, XGBoost vẫn yêu cầu lượng lớn RAM và thời gian tính toán khi xử lý tập dữ liệu có quy mô lớn.
    \item \textbf{Phát triển cây Level-wise:} Chiến lược phát triển theo tầng buộc tất cả các nút ở cùng một độ sâu phải được mở rộng, dẫn đến việc phân chia tại nhiều nút không cần thiết, gây lãng phí tài nguyên tính toán.
\end{itemize}

\paragraph{LightGBM (Light Gradient Boosting Machine)}\hspace{1cm}

LightGBM được thiết kế chuyên biệt để tối ưu hóa hiệu suất trên dữ liệu quy mô lớn thông qua hai kỹ thuật cốt lõi: GOSS (Gradient-based One-Side Sampling) và EFB (Exclusive Feature Bundling). Cây được phát triển theo cơ chế ưu tiên lá (Leaf-wise growth).

Kỹ thuật GOSS ước lượng độ lợi thông tin (Information Gain) bằng cách giữ lại toàn bộ các mẫu có gradient lớn (top $a\%$) và lấy mẫu ngẫu nhiên từ các mẫu có gradient nhỏ (tỷ lệ $b\%$), sau đó nhân trọng số bù:
\begin{equation*}
    \widetilde{V}_j(d) = \frac{1}{n} \left( \frac{\left(\sum_{x_i \in A_l} g_i + \frac{1-a}{b} \sum_{x_i \in B_l} g_i\right)^2}{n_l^j(d)} + \frac{\left(\sum_{x_i \in A_r} g_i + \frac{1-a}{b} \sum_{x_i \in B_r} g_i\right)^2}{n_r^j(d)} \right)
\end{equation*}
trong đó $A$ là tập các mẫu có gradient lớn, $B$ là tập mẫu ngẫu nhiên từ phần gradient nhỏ, $l$ và $r$ ký hiệu nhánh trái và phải tại ngưỡng $d$ của đặc trưng $j$.\\

Kỹ thuật EFB (Exclusive Feature Bundling) gộp các đặc trưng thưa có tính loại trừ lẫn nhau thành các gói (bundles), giảm số lượng đặc trưng hiệu dụng từ $p$ xuống số lượng gói $p' \ll p$.\\

\textbf{Ưu điểm:}
\begin{itemize}
    \item \textbf{Hiệu suất tính toán ưu việt:} Kỹ thuật Histogram giúp giảm độ phức tạp thời gian tìm kiếm điểm chia nhánh xuống $O(\text{số bin} \times \text{số đặc trưng})$, mang lại tốc độ huấn luyện nhanh gấp 3 đến 10 lần so với XGBoost chuẩn.
    \item \textbf{Xử lý biến định danh trực tiếp:} Khai thác hiệu quả các đặc trưng phân loại có lực lượng cao (high-cardinality) như \texttt{building\_id} hay \texttt{meter} bằng cách phân chia tối ưu trên các danh mục mà không cần mã hóa One-Hot, giúp kiểm soát chiều dữ liệu.
    \item \textbf{Tiết kiệm bộ nhớ:} Kỹ thuật EFB giảm đáng kể số chiều hiệu dụng, trong khi biểu diễn Histogram chỉ lưu trữ các bin rời rạc thay vì toàn bộ giá trị liên tục, giúp tiết kiệm RAM khi xử lý dữ liệu hàng chục triệu dòng.
    \item \textbf{GOSS giữ độ chính xác khi lấy mẫu:} Thay vì lấy mẫu đồng đều, GOSS ưu tiên các mẫu có gradient lớn (đóng góp nhiều nhất vào việc học), nên giảm dữ liệu mà vẫn duy trì chất lượng mô hình.
\end{itemize}

\textbf{Nhược điểm:}
\begin{itemize}
    \item \textbf{Rủi ro quá khớp cục bộ:} Sự phát triển theo chiều sâu lá (Leaf-wise) có thể dẫn đến việc mô hình học các mẫu nhiễu cục bộ. Cần giới hạn nghiêm ngặt các tham số \texttt{max\_depth} và \texttt{min\_data\_in\_leaf}.
    \item \textbf{Nhạy cảm với dữ liệu ít:} Chiến lược Leaf-wise tạo ra các cây sâu không cân bằng, dễ overfitting trên các tập dữ liệu nhỏ. Tuy nhiên, hạn chế này ít ảnh hưởng trong bài toán hiện tại do quy mô dữ liệu rất lớn.
\end{itemize}

\paragraph{CatBoost (Categorical Boosting)}\hspace{1cm}

Thuật toán xây dựng các cây quyết định đối xứng (Oblivious Trees) và sử dụng cơ chế Ordered Target Encoding. Đặc điểm cốt lõi của CatBoost là sử dụng hoán vị ngẫu nhiên (random permutation) $\sigma$ để tính giá trị mã hóa cho biến định danh, tránh rò rỉ thông tin mục tiêu:
\begin{equation*}
    \hat{x}_{\sigma_p}^i = \frac{\sum_{j=1}^{p-1} [x_{\sigma_j}^i = x_{\sigma_p}^i] \cdot y_{\sigma_j} + a \cdot P}{\sum_{j=1}^{p-1} [x_{\sigma_j}^i = x_{\sigma_p}^i] + a}
\end{equation*}
trong đó $\sigma_p$ là chỉ số của mẫu thứ $p$ trong hoán vị, $x^i$ là giá trị đặc trưng định danh thứ $i$, $y$ là nhãn mục tiêu, $P$ là prior (thường là trung bình toàn cục của $y$) và $a > 0$ là tham số Laplace smoothing.

Cây đối xứng (Oblivious Tree) sử dụng cùng một điều kiện chia tại tất cả các nút ở cùng một độ sâu, tạo ra cấu trúc cây cân bằng hoàn hảo với $2^d$ lá (với $d$ là chiều sâu). Dự đoán trở thành phép tra bảng $O(d)$.\\

\textbf{Ưu điểm:}
\begin{itemize}
    \item \textbf{Ngăn chặn rò rỉ dữ liệu (Target Leakage):} Ordered Target Encoding tính toán giá trị trung bình mục tiêu dựa trên các mẫu quan sát trong quá khứ nhân tạo theo hoán vị $\sigma$, giúp mã hóa các biến định danh một cách an toàn và mạnh mẽ mà không gây rò rỉ thông tin nhãn.
    \item \textbf{Hiệu năng suy luận (Inference Time):} Cấu trúc cây đối xứng cho phép mô hình được biểu diễn dưới dạng các bảng tra cứu với độ phức tạp suy luận $O(d)$, tăng tốc độ dự đoán lên nhiều lần so với các cây không đối xứng.
    \item \textbf{Ordered Boosting:} Sử dụng nhiều hoán vị khác nhau trong quá trình huấn luyện, giảm thiên kiến (bias) do rò rỉ thông tin trong Gradient Boosting truyền thống, đặc biệt hiệu quả trên dữ liệu có nhiều biến phân loại như \texttt{building\_id}, \texttt{primary\_use}.
    \item \textbf{Kháng overfitting tự nhiên:} Cấu trúc cây đối xứng có ít tham số hơn cây tự do (chỉ $d$ điều kiện chia thay vì $2^d - 1$ nút nội), đóng vai trò như một dạng điều chuẩn ngầm.
\end{itemize}

\textbf{Nhược điểm:}
\begin{itemize}
    \item \textbf{Thời gian huấn luyện trên dữ liệu số thực:} Quy trình mã hóa Ordered Target Encoding và xây dựng cây đối xứng khiến CatBoost mất nhiều thời gian huấn luyện hơn so với LightGBM, đặc biệt khi bộ dữ liệu chứa lượng lớn các biến liên tục.
    \item \textbf{Giới hạn biểu diễn của cây đối xứng:} Ràng buộc cùng một điều kiện chia tại mỗi tầng có thể hạn chế khả năng biểu diễn các mẫu phức tạp, đòi hỏi cây phải sâu hơn để đạt cùng mức chính xác so với cây tự do.
\end{itemize}

\paragraph{Random Forest}\hspace{1cm}

Random Forest là thuật toán học tập hợp Bagging (Bootstrap Aggregating), xây dựng $B$ cây quyết định độc lập trên các tập dữ liệu bootstrap. Mỗi cây được huấn luyện trên tập $\mathcal{D}_b$ được lấy mẫu ngẫu nhiên có hoàn lại (with replacement) từ tập gốc $\mathcal{D}$, với $|\mathcal{D}_b| = |\mathcal{D}| = n$.

Tại mỗi nút phân chia, thay vì xét toàn bộ $p$ đặc trưng, Random Forest chỉ chọn ngẫu nhiên một tập con $m$ đặc trưng và tìm điểm chia tối ưu theo tiêu chí Gini Impurity:
\begin{equation*}
    \text{Gini}(S) = 1 - \sum_{k=1}^{C} p_k^2
\end{equation*}
trong đó $C$ là số lớp và $p_k$ là tỷ lệ mẫu thuộc lớp $k$ trong tập $S$. Số đặc trưng được chọn ngẫu nhiên thường là $m = \lfloor\sqrt{p}\rfloor$ cho bài toán phân lớp.

Dự đoán cuối cùng được xác định bằng bỏ phiếu đa số (majority voting):
\begin{equation*}
    \hat{y}(\mathbf{x}) = \text{mode}\left\{ h_b(\mathbf{x}) \right\}_{b=1}^{B}
\end{equation*}

Phương sai của mô hình tập hợp được giới hạn bởi:
\begin{equation*}
    \text{Var}\left[\hat{y}\right] = \rho \sigma^2 + \frac{1 - \rho}{B} \sigma^2
\end{equation*}
trong đó $\rho$ là hệ số tương quan trung bình giữa các cây và $\sigma^2$ là phương sai của mỗi cây đơn lẻ. Việc chọn ngẫu nhiên $m \ll p$ đặc trưng tại mỗi nút giúp giảm $\rho$, từ đó giảm phương sai tổng thể.\\

\textbf{Ưu điểm:}
\begin{itemize}
    \item \textbf{Độ bền vững với ngoại lai (Robustness):} Cơ chế trung bình hóa từ hàng trăm cây quyết định độc lập giúp mô hình miễn nhiễm với các giá trị bất thường ngẫu nhiên (random spikes) của cảm biến.
    \item \textbf{Huấn luyện song song:} Các cây hoàn toàn độc lập nên có thể huấn luyện song song trên nhiều lõi CPU, tận dụng hiệu quả kiến trúc đa nhân.
    \item \textbf{Ước lượng lỗi ngoài túi (Out-of-Bag Error):} Khoảng $37\%$ mẫu không xuất hiện trong mỗi bootstrap có thể được dùng làm tập kiểm tra nội bộ, cung cấp ước lượng lỗi tổng quát hóa mà không cần tập validation riêng.
    \item \textbf{Đo lường tầm quan trọng đặc trưng (Feature Importance):} Cung cấp thước đo dựa trên mức giảm tạp chất trung bình (Mean Decrease in Impurity) hoặc hoán vị (Permutation Importance), hỗ trợ giải thích mô hình.
\end{itemize}

\textbf{Nhược điểm:}
\begin{itemize}
    \item \textbf{Độ nhạy thấp với lớp thiểu số:} Thuật toán Bagging hội tụ về phía dự đoán nhãn đa số do mỗi bootstrap chỉ chứa một tỷ lệ nhỏ mẫu lớp thiểu số. Cần kết hợp với kỹ thuật Balanced Random Forest (lấy mẫu phân tầng trong bootstrap) để cải thiện.
    \item \textbf{Kích thước mô hình lớn:} Với $B$ cây sâu không giới hạn, file mô hình có thể lên tới hàng GB trên ổ cứng, gây khó khăn cho việc triển khai và lưu trữ.
    \item \textbf{Không có cơ chế học tuần tự:} Khác với Boosting, mỗi cây không học từ sai sót của cây trước, nên không tập trung sửa lỗi ở vùng dữ liệu khó — điểm yếu quan trọng trong bài toán mất cân bằng.
\end{itemize}

\paragraph{Extra Trees (Extremely Randomized Trees)}\hspace{1cm}

Extra Trees (Extremely Randomized Trees) là một biến thể cực đoan hóa của Random Forest. Khác biệt cốt lõi nằm ở hai điểm: Sử dụng toàn bộ tập huấn luyện $\mathcal{D}$ thay vì lấy mẫu bootstrap và ngẫu nhiên hóa hoàn toàn cả ngưỡng phân chia tại mỗi nút.

Cụ thể, tại mỗi nút, thuật toán chọn ngẫu nhiên $K$ đặc trưng. Với mỗi đặc trưng $j$ được chọn, ngưỡng chia $s_j$ được lấy ngẫu nhiên đồng đều trong khoảng giá trị:
\begin{equation*}
    s_j \sim \mathcal{U}\left(\min_{i} x_i^j, \; \max_{i} x_i^j\right)
\end{equation*}
Sau đó, cặp $(j^*, s^*)$ cho điểm chia tốt nhất được chọn theo tiêu chí Gini hoặc Entropy:
\begin{equation*}
    (j^*, s^*) = \arg\max_{j \in \{1,\ldots,K\}, \; s_j} \; \text{Gain}(j, s_j)
\end{equation*}

Việc ngẫu nhiên hóa ngưỡng chia khiến hệ số tương quan $\rho$ giữa các cây giảm sâu hơn so với Random Forest, dẫn đến phương sai tổng thể thấp hơn theo công thức $\text{Var}[\hat{y}] = \rho \sigma^2 + \frac{1-\rho}{B}\sigma^2$.\\

\textbf{Ưu điểm:}
\begin{itemize}
    \item \textbf{Tốc độ huấn luyện nhanh hơn Random Forest:} Không cần tìm kiếm ngưỡng chia tối ưu cho mỗi đặc trưng (chỉ chọn ngẫu nhiên), giúp giảm đáng kể chi phí tính toán tại mỗi nút từ $O(n \log n)$ xuống $O(n)$.
    \item \textbf{Giảm phương sai triệt để hơn:} Hệ số tương quan $\rho$ giữa các cây thấp hơn Random Forest nhờ mức độ ngẫu nhiên hóa cao hơn, giúp mô hình tập hợp có phương sai tổng thể thấp hơn.
    \item \textbf{Kháng overfitting mạnh:} Ngưỡng chia ngẫu nhiên ngăn cây học quá sát các mẫu nhiễu cục bộ. Sử dụng toàn bộ dữ liệu (không bootstrap) giúp mỗi cây có cái nhìn đầy đủ hơn về phân phối dữ liệu, giảm thiên kiến (bias).
    \item \textbf{Phù hợp với dữ liệu nhiễu cao:} Đặc tính ngẫu nhiên hóa cực đoan giúp Extra Trees miễn nhiễm tốt với nhiễu cảm biến, phù hợp với bản chất dữ liệu đo lường năng lượng.
\end{itemize}

\textbf{Nhược điểm:}
\begin{itemize}
    \item \textbf{Thiên kiến (Bias) cao hơn Random Forest:} Ngưỡng chia ngẫu nhiên thay vì tối ưu có thể tạo ra các phân vùng kém chính xác hơn, khiến mỗi cây đơn lẻ có sai số thiên kiến cao hơn.
    \item \textbf{Nhạy cảm với đặc trưng không liên quan:} Do không tìm kiếm ngưỡng tối ưu, thuật toán khó phân biệt giữa các đặc trưng hữu ích và đặc trưng nhiễu, đặc biệt khi tập đặc trưng có chiều cao.
    \item \textbf{Cùng hạn chế Bagging với lớp thiểu số:} Tương tự Random Forest, Extra Trees không có cơ chế tập trung vào các mẫu khó (hard samples) của lớp thiểu số, kém hiệu quả hơn Boosting trong bài toán mất cân bằng nghiêm trọng.
\end{itemize}
